\section{Adaptation lineage and calibration controls}
\label{app:adaptation-lineage}
\subsection{Adaptation lineage and rejected candidates}
The name ``adapted PHASE'' refers to several checkpoints with different fitting populations. F0 denotes the original nine-component TRAIN40 graph. The first F1 adaptation updates its original 1,955-parameter fusion head while freezing the eight preceding components and their normalizers. Its all-32-recording refit uses 88,744 TRAIN40 OOF states and 3,571 admitted EvTTC states, for 92,315 training rows; 257 native queries lack compatible targets. Each of seeds 7, 13, and 21 completes two epochs and 1,444 updates, with seed 7 designated for delivery before scoring. These are training-fit results. \evidence{A01}

The preserved all-32 F1 cache confirms that its supervision consumed the old CCRs1-low/medium source assignment. Its preparation script reads each query's \texttt{ttc\_csv\_path}, hashes the file, and joins its TTC curve before writing supervised shards. The preserved \texttt{data\_v3/TARGET\_LEDGER.json} records the old low-file hash beginning \texttt{780ee30f} and medium-file hash beginning \texttt{1c5adba8}; under the corrected identity map those files belong to medium and low, respectively. It admits 98 low and 141 medium queries, with two and eight additional unsupported targets. Its completion receipt binds both that ledger and the merged training array, and the three seeds' start receipts and delivered seed-7 bundle bind the same completion receipt. Thus 239 admitted supervised rows used the old map, and the ten exclusions also require a new eligibility check. Rescoring retained outputs cannot repair those F1 weights: a corrected all-32 F1 requires rebuilding the supervised cache and refitting seeds 7/13/21. No corrected all-32 F1 training is reported here. The grouped F1 campaign below excludes the eight overlapping competition recordings, including these two recordings, and restarts from F0 in each fold. This particular source swap does not invalidate that grouped comparison or establish a dependency in the final 688-parameter router, whose optimizer population excludes those recordings. \evidence{A01,A02}

That grouped campaign comprises 21 head fits: seven scenario families times three seeds, with 2,531 queries from twenty-four road recordings. It retains TRAIN40 during each fit. F1 reduces recording-macro RAE from F0's 30.600\% to 25.468/26.188/25.938\%, but increases query MAE from 0.8494\,s to 0.8585/0.8925/0.8653\,s for seeds 7/13/21. All three family-bootstrap intervals for the RAE difference include zero. For targets above 4\,s, seed 7 increases MAE from 1.2542 to 1.5496\,s. On Slider-500, outside this optimizer fit but historically exposed, F1 increases RAE from 83.887\% to 87.578/91.022/87.330\%; each variant returns 76 of 77 outputs. Lower road relative error therefore did not establish uniform improvement. \evidence{A02}

The next quality-fusion campaign freezes all nine F0 components and fits three 401-parameter gates. Its named \texttt{evaluated} checkpoint has separate fitting, calibration, and confirmation roles; \texttt{all\_data\_refit} instead consumes all 2,677 eligible queries from twenty-four road and three Slider recordings. The latter's reported errors are resubstitution. On 1,747 road OOF queries, the quality gate gives RAE 12.834\% and MAE 0.4224\,s, versus 22.081\% and 0.7544\,s for an equal-parameter control without quality descriptors. Both receive geometry. This supports quality information for that gate and population. \evidence{A03}

The preserved bundles close the upstream-weight distinction: the 401-parameter gate, startup router, and submitted router all bind the original F0 head. Hash checks of their thirteen upstream components, including normalizers and seal/recipe records, agree with their manifests. None substitutes the affected all-32 F1 heads; complete checkpoint identities and cache bindings are retained in the additional source index. The gate preparation explicitly selects F0, RGB context, and event context as native experts; startup reuses those experts, and submission refits the router while copying the verified original upstream bundle. This establishes a separate weight lineage, while retaining the historical data exposure of all three gate campaigns. \evidence{A01,A03,A04,A07}

The startup campaign then executes 117 fits, including Q/QA/NAcap controls, seconds- and inverse-coordinate trials, and all-data refits. It retains the 688-parameter inverse-coordinate router as a candidate after rejecting the seconds-coordinate trial. Road OOF MAE changes from the previous gate's 0.4224 to 0.4049\,s, while CCRm confirmation MAE worsens from 1.020 to 1.036\,s. Its descriptive road RAE-difference interval, resampling five families, is $[-1.5176,0.8001]$ percentage points. Development now includes Slider-750/1000, so this comparison changes both architecture and exposure. The final submitted recipe below returns to a 2,608-query optimizer population; earlier use of those Slider recordings still counts as design exposure. \evidence{A04}

A further 312 completed small-head fits test sign/log/geometry variants (96), supervised cost (48), correlated expert perturbations (24), native-context resets (48), support versus capacity (48), and support with resets (48). Eight family-held-out folds, seeds 7/13/21, and 300 updates per fit use the 2,608-query population. Promotion requires full coverage and no worsening of both recording-macro RAE and MAE in either road or Slider relative to IQ. None passes all four error criteria. The six preserved fits from an earlier failed launch do not enter the 312 total. The 447 context resets cover 149 prepared queries; they do not represent 447 cold sensor starts. These fits are separate from the 21 F1 fits, the 117 startup fits, and the later 48 canonical-router fits. \evidence{A05}

\subsection{Measured interval coverage by variant}
The grouped F1 experiment also measures substantial undercoverage. Its nominal 80\% intervals cover 67.88/65.78/67.72\% for seeds 7/13/21. Rescaling with each fold's training residuals raises coverage only to 70.17/68.63/70.01\%; that procedure is not split conformal and was not bundled. The earlier 401-parameter \texttt{evaluated} gate uses separate CPNAO calibration records and attains 79.83\% CCRm coverage. The startup \texttt{evaluated} candidate reports 83.10\% CCRm coverage after separate CPNAO calibration. These are descriptive results on historically exposed groups. Their all-data refits consume those calibration labels and carry no validated equivalent interval; neither does the final 2,608-query submitted router. Finite-output coverage and probabilistic interval coverage describe different properties. \evidence{A02,A03,A04,A07}


\subsection{Geometry controls and support differences}
\label{app:geometry-attribution}
F0 versus the final router changes both the available experts and the learned combination. The older 401-parameter capacity control establishes a benefit of quality descriptors only within its own experiment. Its Slider-750 geometry observer has MAE 0.1976\,s on 49/51 outputs, versus 3.3157\,s for that gate on 51/51; Slider-1000 gives 0.0374\,s on 16/18 versus 0.2757\,s on 18/18. Different masks prevent a direct attribution from these aggregates. The archived corrected 434-query evaluator does not provide a complete matched geometry-only, geometry-with-F0-fallback, fixed-mixture, and final-router comparison with common support and timing. A matched comparison is needed to separate the contribution of geometry from that of learned routing; the present gain belongs to the combined adapter. \evidence{A03,N01}
